\documentclass[11pt,a4paper]{article}
\usepackage[T1]{fontenc}\usepackage[utf8]{inputenc}\usepackage[provide=*,italian]{babel}
\usepackage{lmodern,microtype,geometry}\geometry{margin=2.5cm}
\usepackage{amsmath,amssymb,amsthm,mathtools,booktabs,array,longtable,xcolor,enumitem,tikz}
\usetikzlibrary{arrows.meta,positioning,calc}
\usepackage[most]{tcolorbox}\usepackage{listings,hyperref,fancyhdr}
\definecolor{sapienzared}{RGB}{130,0,36}\definecolor{softred}{RGB}{249,241,243}
\definecolor{softblue}{RGB}{239,246,252}\definecolor{softgray}{RGB}{245,245,245}\definecolor{codegreen}{RGB}{25,110,65}
\hypersetup{colorlinks=true,linkcolor=sapienzared,urlcolor=sapienzared,pdftitle={Tecniche reticolari e CPM},pdfauthor={Giampaolo Liuzzi}}
\pagestyle{fancy}\fancyhf{}\fancyhead[L]{\small Ottimizzazione su Reti}\fancyhead[R]{\small Capitolo 8}\fancyfoot[C]{\thepage}\renewcommand{\headrulewidth}{0.4pt}
\newtheorem{definition}{Definizione}[section]\newtheorem{theorem}[definition]{Teorema}\newtheorem{proposition}[definition]{Proposizione}
\newtcolorbox{learningbox}{colback=softred,colframe=sapienzared,title=Obiettivi di apprendimento,fonttitle=\bfseries,boxrule=0.8pt,arc=1.5mm}
\newtcolorbox{examplebox}[1]{colback=softblue,colframe=blue!55!black,title=#1,fonttitle=\bfseries,boxrule=0.7pt,arc=1.5mm}
\newtcolorbox{notationbox}{colback=softgray,colframe=black!55,title=Registro delle notazioni,fonttitle=\bfseries,boxrule=0.7pt,arc=1.5mm}
\newtcolorbox{warningbox}{colback=yellow!9,colframe=orange!65!black,title=Attenzione,fonttitle=\bfseries,boxrule=0.7pt,arc=1.5mm}
\newtcolorbox{networkxbox}{colback=softgray,colframe=black!55,title=Python/NetworkX,fonttitle=\bfseries,boxrule=0.7pt,arc=1.5mm}
\lstdefinestyle{python}{language=Python,basicstyle=\ttfamily\small,keywordstyle=\color{blue!65!black}\bfseries,commentstyle=\color{codegreen},stringstyle=\color{sapienzared},backgroundcolor=\color{softgray},frame=single,rulecolor=\color{black!25},showstringspaces=false,breaklines=true,columns=fullflexible}
\newcommand{\V}{\mathcal{V}}\newcommand{\E}{\mathcal{E}}
\begin{document}\hypersetup{pageanchor=false}
\begin{titlepage}\centering\vspace*{1.6cm}
{\color{sapienzared}\rule{\textwidth}{1.2pt}}\\[1.1cm]{\Large Ottimizzazione su Reti}\\[0.35cm]
{\huge\bfseries Capitolo 8}\\[0.45cm]{\LARGE\bfseries Tecniche reticolari e CPM}\\[1.1cm]
{\color{sapienzared}\rule{0.72\textwidth}{0.8pt}}\\[1.1cm]
{\large Corso di Laurea Magistrale in Ingegneria Gestionale}\\[0.25cm]{\large Sapienza Università di Roma}\\[1.2cm]
{\large Giampaolo Liuzzi}\\[0.25cm]{\normalsize Dipartimento di Ingegneria Informatica, Automatica e Gestionale}
\vfill{\small Versione preliminare per uso didattico}\end{titlepage}
\hypersetup{pageanchor=true}\pagenumbering{roman}\tableofcontents\newpage\pagenumbering{arabic}

\begin{learningbox}
Al termine del capitolo lo studente dovrebbe essere in grado di:
\begin{itemize}[nosep]
\item costruire una rete di progetto a partire dalle precedenze;
\item distinguere rappresentazioni attività-sugli-archi e attività-sui-nodi;
\item calcolare tempi al più presto e al più tardi;
\item determinare slittamenti, attività critiche e cammini critici;
\item interpretare il CPM come cammino massimo su un DAG;
\item effettuare analisi di sensitività e implementare il calcolo in NetworkX.
\end{itemize}
\end{learningbox}

\section{Progetti e relazioni di precedenza}
Un progetto è composto da attività di durata deterministica. Le relazioni di precedenza specificano quali attività devono essere completate prima che altre possano iniziare. Il \emph{Critical Path Method} (CPM) determina la durata minima del progetto, le date al più presto e al più tardi, le attività critiche e uno o più cammini critici.

Le precedenze devono definire un DAG. Un ciclo significherebbe che un'attività dipende, direttamente o indirettamente, da se stessa e renderebbe il progetto logicamente irrealizzabile.

\section{Due rappresentazioni della rete}
\subsection{Attività sugli archi}
Nella rappresentazione AOA (\emph{activity on arc}) i nodi sono eventi e ogni attività è un arco $(i,j)$ di durata $p_{ij}$. L'evento $i$ deve verificarsi prima che l'attività inizi; l'evento $j$ si verifica dopo il completamento di tutte le attività entranti.

Talvolta sono necessarie \emph{attività fittizie} di durata zero per descrivere correttamente le precedenze. La rete viene completata con un unico evento iniziale $0$ e un unico evento finale $f$.

\subsection{Attività sui nodi}
Nella rappresentazione AON (\emph{activity on node}) ogni attività è un nodo e un arco $a\to b$ indica che $a$ precede immediatamente $b$. Le durate sono attributi dei nodi.
\begin{center}\begin{tabular}{lll}\toprule Rappresentazione&Attività&Precedenze\\\midrule AOA&archi&eventi e archi anche fittizi\\ AON&nodi&archi del DAG\\\bottomrule\end{tabular}\end{center}
Nel seguito adottiamo principalmente AOA, coerentemente con la formulazione del corso.

\section{CPM come cammino massimo in un DAG}
In una rete AOA, un evento può verificarsi soltanto quando tutte le attività entranti sono terminate. Se $t_i$ è il tempo al più presto dell'evento $i$, allora
\begin{equation}\label{eq:forward}
t_0=0,\qquad t_j=\max_{(i,j)\in\E}\{t_i+p_{ij}\}.
\end{equation}
Sono le equazioni di Bellman per un cammino \emph{massimo}. Poiché la rete è aciclica, le etichette si calcolano una sola volta in ordine topologico. La durata minima del progetto è
\[T=t_f,\]
ossia la lunghezza del cammino massimo dall'evento iniziale a quello finale.

\begin{examplebox}{Dai cammini minimi ai cammini massimi}
Si può porre $c_{ij}=-p_{ij}$ e calcolare un cammino minimo nel DAG. Nel CPM è però più naturale mantenere durate positive e usare direttamente il massimo nella ricorrenza \eqref{eq:forward}.
\end{examplebox}

\section{Scansione in avanti: tempi al più presto}
\begin{definition}[Tempo al più presto]
Il tempo $t_i$ è il primo istante in cui può verificarsi l'evento $i$, cioè il primo istante entro cui possono essere completate tutte le attività afferenti a $i$.
\end{definition}
Dato un ordinamento topologico $0,1,\ldots,f$, si pone $t_0=0$ e si applica \eqref{eq:forward}. Per l'attività $(i,j)$:
\[ES_{ij}=t_i,\qquad EF_{ij}=t_i+p_{ij},\]
dove $ES$ ed $EF$ sono inizio e fine al più presto. Il tempo minimo di completamento è $C_{ij}=EF_{ij}$.

\section{Scansione all'indietro: tempi al più tardi}
\begin{definition}[Tempo al più tardi]
Il tempo $T_i$ è l'ultimo istante in cui può verificarsi l'evento $i$ senza aumentare la durata minima del progetto.
\end{definition}
Si pone $T_f=t_f$ e si percorrono i nodi in ordine topologico inverso:
\begin{equation}\label{eq:backward}
T_i=\min_{(i,j)\in\E}\{T_j-p_{ij}\}.
\end{equation}
Per l'attività $(i,j)$:
\[LF_{ij}=T_j,\qquad LS_{ij}=T_j-p_{ij}.\]
In generale $t_i\leq T_i$; $T_i-t_i$ è il margine dell'evento.

\section{Slittamenti e attività critiche}
\begin{definition}[Slittamento totale]
Lo slittamento totale dell'attività $(i,j)$ è
\begin{equation}\label{eq:slack}
TS_{ij}=LS_{ij}-ES_{ij}=T_j-t_i-p_{ij}.
\end{equation}
\end{definition}
Si ha sempre $TS_{ij}\geq0$. Un ritardo non superiore a $TS_{ij}$ può essere assorbito senza posticipare il progetto, a parità del resto.

\begin{definition}[Attività critica]
Un'attività è critica se $TS_{ij}=0$. Un suo ritardo positivo aumenta la durata minima del progetto dello stesso ammontare, a parità delle altre durate.
\end{definition}

Lo slittamento libero è
\[FS_{ij}=t_j-t_i-p_{ij}.\]
Esso misura il ritardo ammissibile senza posticipare il tempo al più presto di $j$, e vale $0\leq FS_{ij}\leq TS_{ij}$.

\section{Cammini critici}
\begin{definition}[Cammino critico]
Un cammino dall'evento iniziale a quello finale è critico se la somma delle durate delle sue attività è uguale a $T$.
\end{definition}
Ogni attività di un cammino critico è critica. Viceversa, l'insieme delle attività critiche può contenere più cammini e ramificazioni: occorre verificarne la continuità da $0$ a $f$.

\begin{proposition}
Un arco $(i,j)$ appartiene a un cammino critico se
\[t_i=T_i,\qquad t_j=T_j,\qquad t_j=t_i+p_{ij}.\]
\end{proposition}

\begin{warningbox}
Un progetto può avere più cammini critici. Ridurre la durata di una sola attività critica può allora non ridurre la durata del progetto: occorre intervenire su tutti i cammini che restano massimi.
\end{warningbox}

\section{Esempio completo}
Consideriamo la rete AOA seguente.
\begin{center}\small
\begin{tabular}{cccccccccc}\toprule
Attività&A&B&C&D&E&F&G&H&I\\
Arco&$(0,1)$&$(1,2)$&$(1,3)$&$(2,4)$&$(2,5)$&$(3,5)$&$(4,6)$&$(5,6)$&$(6,7)$\\
Durata&5&7&3&4&5&6&2&8&2\\\bottomrule
\end{tabular}\end{center}

\subsection{Passaggio in avanti}
\[\begin{array}{c|rrrrrrrr}i&0&1&2&3&4&5&6&7\\\hline t_i&0&5&12&8&16&17&25&27\end{array}\]
Per esempio,
\[t_5=\max\{t_2+5,t_3+6\}=\max\{17,14\}=17,\]
e $T=t_7=27$.

\subsection{Passaggio all'indietro}
Ponendo $T_7=27$ si ottiene
\[\begin{array}{c|rrrrrrrr}i&0&1&2&3&4&5&6&7\\\hline T_i&0&5&12&11&23&17&25&27\end{array}\]
Per esempio,
\[T_2=\min\{T_4-4,T_5-5\}=\min\{19,12\}=12.\]

\subsection{Slittamenti}
\begin{center}\begin{tabular}{crrrrrrrrr}\toprule Attività&A&B&C&D&E&F&G&H&I\\\midrule $TS$&0&0&3&7&0&3&7&0&0\\\bottomrule\end{tabular}\end{center}
Le attività critiche sono A, B, E, H e I; il cammino critico è
\[0\xrightarrow{A}1\xrightarrow{B}2\xrightarrow{E}5\xrightarrow{H}6\xrightarrow{I}7,\]
di durata $5+7+5+8+2=27$.

\section{Rappresentazione AON}
In una rete AON, per ogni attività $a$ di durata $p_a$:
\[ES_a=\max_{h\in\operatorname{Pred}(a)}EF_h,\qquad EF_a=ES_a+p_a.\]
Se $T=\max_a EF_a$, il passaggio all'indietro usa
\[LF_a=\min_{k\in\operatorname{Succ}(a)}LS_k,\qquad LS_a=LF_a-p_a,\]
ponendo $LF_a=T$ per le attività terminali. Lo slittamento è $LS_a-ES_a$. È importante non mescolare formule AOA e AON senza dichiarare la rappresentazione.

\section{Sensitività delle durate}
Se la durata di un'attività non critica aumenta di una quantità non superiore al suo slittamento totale, $T$ non cambia. Superato il margine può emergere un nuovo cammino critico.

Per ridurre $T$ si può accelerare, o \emph{crashare}, un'attività critica sostenendo un costo addizionale. Dopo ogni modifica il CPM va ricalcolato, perché il cammino critico può cambiare o moltiplicarsi.

\begin{examplebox}{Interpretazione del margine}
Nell'esempio, D ha $TS=7$. Un aumento della sua durata da $4$ a $10$ lascia invariato $T=27$; portandola a $12$ si supera il margine di un giorno e il percorso attraverso D e G diventa più lungo del precedente cammino critico.
\end{examplebox}

\section{Python e NetworkX}
\begin{networkxbox}
NetworkX non espone una funzione CPM completa, ma offre ordinamento topologico e cammino massimo nei DAG. Una breve implementazione produce tempi al più presto, al più tardi e slittamenti.
\end{networkxbox}

\begin{lstlisting}[style=python,caption={CPM per una rete AOA}]
import networkx as nx

G = nx.DiGraph()
activities = [
 (0,1,"A",5), (1,2,"B",7), (1,3,"C",3),
 (2,4,"D",4), (2,5,"E",5), (3,5,"F",6),
 (4,6,"G",2), (5,6,"H",8), (6,7,"I",2)
]
for i, j, name, duration in activities:
    G.add_edge(i, j, name=name, duration=duration)

order = list(nx.topological_sort(G))
t = {v: float("-inf") for v in G}
t[order[0]] = 0
for i in order:
    for j in G.successors(i):
        t[j] = max(t[j], t[i] + G[i][j]["duration"])

finish = order[-1]
T = {v: t[finish] for v in G}
for i in reversed(order[:-1]):
    T[i] = min(T[j] - G[i][j]["duration"]
               for j in G.successors(i))

for i, j, data in G.edges(data=True):
    slack = T[j] - t[i] - data["duration"]
    print(data["name"], slack)
\end{lstlisting}
In generale occorre gestire più nodi iniziali o finali con eventi fittizi, verificare l'aciclicità e confrontare gli slittamenti in virgola mobile con una tolleranza.

\section{Registro delle notazioni}
\begin{notationbox}\begin{center}\begin{tabular}{ll}\toprule
Simbolo&Significato\\\midrule
$p_{ij}$&durata dell'attività $(i,j)$\\
$t_i$&tempo al più presto dell'evento $i$\\
$T_i$&tempo al più tardi dell'evento $i$\\
$T=t_f$&durata minima del progetto\\
$ES,EF$&inizio e fine al più presto\\
$LS,LF$&inizio e fine al più tardi\\
$C_{ij}$&$t_i+p_{ij}$, completamento al più presto\\
$TS_{ij}$&$T_j-t_i-p_{ij}$, slittamento totale\\
$FS_{ij}$&$t_j-t_i-p_{ij}$, slittamento libero\\\bottomrule
\end{tabular}\end{center}\end{notationbox}

\section{Errori frequenti}
\begin{enumerate}
\item Usare il minimo nel passaggio in avanti o il massimo in quello all'indietro.
\item Confondere $t_i$ con la durata dell'attività.
\item Calcolare $TS_{ij}$ usando $T_i$ invece di $T_j$.
\item Dichiarare critica un'attività soltanto perché ha durata elevata.
\item Confondere attività critica e cammino critico.
\item Dimenticare attività fittizie o attribuire loro durata positiva.
\item Mescolare formule AOA e AON.
\item Assumere che ridurre una sola attività critica riduca sempre il progetto.
\end{enumerate}

\section{Esercizi}
\textbf{Esercizio 8.1 -- Costruzione.} Da attività, durate e precedenze costruire sia AON sia una corretta rete AOA.

\medskip\textbf{Esercizio 8.2 -- Passaggio in avanti.} Calcolare ordinamento topologico, $t_i$, $ES$, $EF$ e durata minima.

\medskip\textbf{Esercizio 8.3 -- Passaggio all'indietro.} Calcolare $T_i$, $LS$, $LF$, slittamenti totali e liberi.

\medskip\textbf{Esercizio 8.4 -- Cammini multipli.} Modificare una durata in modo che esistano due cammini critici.

\medskip\textbf{Esercizio 8.5 -- Sensitività.} Per ogni attività non critica determinare il massimo aumento che lascia invariato $T$.

\medskip\textbf{Esercizio 8.6 -- NetworkX.} Implementare l'esempio, verificare la durata con \path{dag_longest_path_length} e produrre una tabella riepilogativa.

\section{Sintesi del capitolo}
\begin{learningbox}\begin{itemize}[nosep]
\item Una rete di progetto è un DAG che rappresenta attività e precedenze.
\item Il passaggio in avanti è un calcolo di cammini massimi in ordine topologico.
\item Il passaggio all'indietro determina gli ultimi tempi compatibili con $T$.
\item Lo slittamento totale misura il ritardo assorbibile da un'attività.
\item Le attività critiche hanno slittamento nullo e formano uno o più cammini critici.
\item Dopo una variazione delle durate il CPM deve essere ricalcolato.
\end{itemize}\end{learningbox}

\paragraph{Conclusione.} Il CPM chiude il percorso mostrando come grafi aciclici, ordinamenti topologici, cammini e modelli di flusso forniscano un linguaggio unitario per problemi apparentemente diversi.
\end{document}
